% ============================================================
%  PRÉSENTATION DE SOUTENANCE — BrainScan AI
%  Auteurs   : MOUMEN Doae Malak & SOUALHINE Douaa
%  Encadrant : Pr. BENLAKHDAR Said
%  École     : EMSI
%  Compiler  : pdflatex  (Overleaf : Ctrl+S → compile)
% ============================================================

\documentclass[aspectratio=169, 12pt, compress]{beamer}

% ── Encodage & langue ──────────────────────────────────────
\usepackage[utf8]{inputenc}
\usepackage[T1]{fontenc}
\usepackage[french]{babel}

% ── Graphiques ─────────────────────────────────────────────
\usepackage{graphicx}
\usepackage{tikz}
\usetikzlibrary{arrows.meta, positioning, shapes.geometric, shadows, overlay-beamer-styles}

% ── Tableaux & maths ───────────────────────────────────────
\usepackage{booktabs}
\usepackage{tabularx}
\usepackage{colortbl}
\usepackage{amsmath}
\usepackage{amssymb}

% ── Icônes légères sans FontAwesome ────────────────────────
\usepackage{pifont}
\newcommand{\cmark}{\textcolor{EMSIgreen}{\ding{51}}}
\newcommand{\xmark}{\textcolor{EMSIred}{\ding{55}}}

% ── Boîtes colorées dans Beamer ────────────────────────────
\usepackage{tcolorbox}
\tcbuselibrary{skins}

% ============================================================
%  THÈME ET COULEURS
% ============================================================
\usetheme{Madrid}
\usecolortheme{default}

% Palette EMSI — bleu marine + gris acier
\definecolor{EMSIblue}{RGB}{12,  55, 103}   % bleu marine EMSI
\definecolor{EMSIblue2}{RGB}{0, 102, 179}   % bleu clair accent
\definecolor{EMSIgray}{RGB}{75,  85,  95}   % gris acier
\definecolor{EMSIlight}{RGB}{235,242,250}   % bleu très clair fond
\definecolor{EMSIgreen}{RGB}{39, 174,  96}  % vert verdict sain
\definecolor{EMSIred}{RGB}{192,  57,  43}   % rouge verdict tumeur
\definecolor{EMSIorange}{RGB}{211,  84,   0} % orange zone incertitude

% Application des couleurs sur la structure Beamer
\setbeamercolor{structure}{fg=EMSIblue}
\setbeamercolor{palette primary}{bg=EMSIblue,    fg=white}
\setbeamercolor{palette secondary}{bg=EMSIblue2, fg=white}
\setbeamercolor{palette tertiary}{bg=EMSIgray,   fg=white}
\setbeamercolor{palette quaternary}{bg=EMSIblue, fg=white}
\setbeamercolor{frametitle}{bg=EMSIblue, fg=white}
\setbeamercolor{title}{bg=EMSIblue, fg=white}
\setbeamercolor{block title}{bg=EMSIblue2, fg=white}
\setbeamercolor{block body}{bg=EMSIlight, fg=black}
\setbeamercolor{block title alerted}{bg=EMSIred, fg=white}
\setbeamercolor{block body alerted}{bg=red!5, fg=black}
\setbeamercolor{block title example}{bg=EMSIgreen!80!black, fg=white}
\setbeamercolor{block body example}{bg=green!5, fg=black}

% Police sans-serif propre
\usefonttheme{professionalfonts}
\setbeamerfont{frametitle}{series=\bfseries, size=\large}
\setbeamerfont{title}{series=\bfseries, size=\LARGE}

% Bullets ronds et lisibles
\setbeamertemplate{itemize item}{\color{EMSIblue2}\textbullet}
\setbeamertemplate{itemize subitem}{\color{EMSIgray}$\circ$}
\setbeamertemplate{navigation symbols}{}   % supprimer les boutons nav

% ============================================================
%  LOGO EMSI — coin supérieur gauche sur chaque diapositive
%  --> Remplacez 'logo_emsi' par le chemin réel de votre logo
%      (ex. : logo_emsi.png placé dans le même dossier)
% ============================================================
\addtobeamertemplate{frametitle}{}{%
  \begin{tikzpicture}[remember picture, overlay]
    \node[anchor=north west, inner sep=5pt]
      at (current page.north west)
      {\includegraphics[height=0.72cm]{logo_emsi}};
  \end{tikzpicture}%
}

% ============================================================
%  MÉTADONNÉES
% ============================================================
\title[BrainScan AI]{%
  Détection Automatique de Tumeurs Cérébrales\\[4pt]
  par Deep Learning\\[2pt]
  {\normalsize\textit{BrainScan AI}}}

\subtitle{Projet de Fin d'Études}

\author[MOUMEN \& SOUALHINE]{%
  \textbf{MOUMEN Doae Malak} \quad\&\quad \textbf{SOUALHINE Douaa}\\[4pt]
  {\small Encadrant : Pr.~BENLAKHDAR Said}}

\institute[EMSI]{%
  École Marocaine des Sciences de l'Ingénieur (EMSI)\\
  Département Informatique \& Intelligence Artificielle}

\date{Soutenance PFE — 2024--2025}

% ============================================================
%  COMMANDES UTILITAIRES
% ============================================================
\newcommand{\code}[1]{\texttt{\small #1}}

% Boîte résultat colorée
\newtcolorbox{resultbox}[2][]{
  enhanced, sharp corners,
  colback=#2!8, colframe=#2!70!black,
  fonttitle=\bfseries\small, title=#1,
  left=4pt, right=4pt, top=3pt, bottom=3pt,
  boxrule=0.6pt,
}

% ============================================================
\begin{document}
% ============================================================

% ── PAGE DE GARDE ──────────────────────────────────────────
\begin{frame}[plain]
  \begin{tikzpicture}[remember picture, overlay]
    % Bandeau supérieur bleu
    \fill[EMSIblue] (current page.north west)
      rectangle ([yshift=-1.8cm]current page.north east);
    % Bandeau inférieur gris
    \fill[EMSIgray!80] (current page.south west)
      rectangle ([yshift=1.0cm]current page.south east);
    % Logo EMSI — haut gauche
    \node[anchor=north west, inner sep=10pt]
      at (current page.north west)
      {\includegraphics[height=1.1cm]{logo_emsi}};
    % Texte institution
    \node[anchor=north east, inner sep=10pt, white, font=\small\itshape]
      at (current page.north east)
      {École Marocaine des Sciences de l'Ingénieur};
  \end{tikzpicture}

  \vspace{0.5cm}
  \titlepage
\end{frame}

% ── TABLE DES MATIÈRES ─────────────────────────────────────
\begin{frame}{Plan de la Présentation}
  \tableofcontents[hideallsubsections]
\end{frame}

% ============================================================
%  SECTION 1 — INTRODUCTION & PROBLÉMATIQUE
% ============================================================
\section{Introduction et Problématique}

\begin{frame}{Contexte Médical et Enjeux}

  \begin{columns}[T]

    % Colonne gauche : chiffres clés
    \begin{column}{0.44\textwidth}
      \begin{alertblock}{Données épidémiologiques}
        \begin{itemize}
          \item Taux de survie à 5 ans\\
                \textbf{< 10\,\%} pour les glioblastomes
          \item Diagnostic manuel~: des centaines de coupes
                IRM par patient, sujet à l'erreur humaine
          \item \textbf{Détection précoce} = principal
                facteur pronostique
        \end{itemize}
      \end{alertblock}

      \vspace{6pt}

      \begin{block}{Verrou scientifique}
        Comment automatiser la lecture d'un volume IRM
        3D pour détecter une lésion en quelques secondes,
        sans expertise radiologique sur place~?
      \end{block}
    \end{column}

    % Colonne droite : proposition de valeur
    \begin{column}{0.52\textwidth}
      \begin{exampleblock}{Notre réponse~: BrainScan AI}
        \begin{itemize}
          \item[\cmark] Application web \textbf{Django}
                        accessible via navigateur
          \item[\cmark] Modèle \textbf{ResNet50} fine-tuné
                        sur BraTS~2021
          \item[\cmark] Résultat en \textbf{< 3 secondes}
                        après upload
          \item[\cmark] Interface orientée
                        \textit{réassurance} patient
          \item[\cmark] Seuil de décision calculé
                        \textbf{automatiquement} (Youden)
        \end{itemize}
      \end{exampleblock}

      \vspace{6pt}
      \begin{tikzpicture}
        \node[rounded corners=4pt, fill=EMSIlight,
              draw=EMSIblue2, inner sep=6pt,
              font=\small, text width=0.90\columnwidth]{%
          \textbf{Dataset~:} BraTS 2021 \\
          \textbf{Format entrée~:} \code{.nii.gz} (NIfTI) \\
          \textbf{Sortie~:} Probabilité tumorale + verdict
        };
      \end{tikzpicture}
    \end{column}

  \end{columns}
\end{frame}

% ============================================================
%  SECTION 2 — MÉTHODOLOGIE TECHNIQUE
% ============================================================
\section{Méthodologie Technique}

\begin{frame}{Architecture et Pipeline de Traitement}

  % Diagramme pipeline (TikZ)
  \begin{center}
  \begin{tikzpicture}[
    node distance=0.35cm and 0.55cm,
    box/.style={
      draw=EMSIblue, fill=EMSIlight,
      rounded corners=3pt, align=center,
      minimum height=0.72cm, minimum width=2.1cm,
      font=\scriptsize\bfseries
    },
    arr/.style={-{Stealth[length=4pt]}, EMSIblue2, thick}
  ]
    \node[box] (A) {IRM\\NIfTI};
    \node[box, right=of A] (B) {Coupe\\axiale \#75};
    \node[box, right=of B] (C) {Z-Score\\Dataset};
    \node[box, right=of C] (D) {CLAHE};
    \node[box, right=of D] (E) {ResNet50\\224×224};
    \node[box, right=of E] (F) {Inversion\\$1-p$};
    \node[box, right=of F] (G) {Seuil\\Youden};

    \draw[arr] (A)--(B); \draw[arr] (B)--(C);
    \draw[arr] (C)--(D); \draw[arr] (D)--(E);
    \draw[arr] (E)--(F); \draw[arr] (F)--(G);
  \end{tikzpicture}
  \end{center}

  \vspace{4pt}

  \begin{columns}[T]

    \begin{column}{0.31\textwidth}
      \begin{block}{ResNet50 — Fine-Tuning}
        \begin{itemize}\small
          \item Poids ImageNet
          \item 20 dernières couches entraînables
          \item L2 + Dropout~0.5
          \item Sortie sigmoïde
        \end{itemize}
      \end{block}
    \end{column}

    \begin{column}{0.33\textwidth}
      \begin{block}{Z-Score Dataset-Level}
        \[
          I' = \frac{I - \mu_{\text{train}}}{\sigma_{\text{train}}}
        \]
        \small Statistiques calculées une seule fois,
        sauvegardées dans \code{dataset\_stats.json}
      \end{block}
    \end{column}

    \begin{column}{0.32\textwidth}
      \begin{block}{Focal Loss + Label Smoothing}
        \[
          \mathcal{L} = \mathbb{E}[w_f]\cdot\mathcal{L}_{\text{BCE}}^{\varepsilon}
        \]
        \small $\gamma=2$, $\alpha=0.25$, $\varepsilon=0.1$\\
        Contre la surconfiance médicale
      \end{block}
    \end{column}

  \end{columns}

  \vspace{4pt}
  \begin{center}
  \begin{tikzpicture}
    \node[fill=EMSIred!10, draw=EMSIred!60, rounded corners=3pt,
          inner sep=5pt, font=\small]{%
      \textbf{Anomalie corrigée~:} classes inversées à l'entraînement
      (0\,=\,Tumeur, 1\,=\,Sain)
      $\;\Rightarrow\;p_{\text{tumeur}} = 1 - p_{\text{brute}}$
    };
  \end{tikzpicture}
  \end{center}

\end{frame}

% ============================================================
%  SECTION 3 — DÉMONSTRATION INTERFACE
% ============================================================
\section{Démonstration Interface}

% ── Slide 3 : Résultat Sain ────────────────────────────────
\begin{frame}{Démonstration — Diagnostic~: Cerveau Sain}

  \begin{columns}[T]

    % Capture d'écran
    \begin{column}{0.55\textwidth}
      \begin{tikzpicture}
        \node[draw=EMSIgreen!70, line width=1.2pt,
              rounded corners=4pt, inner sep=3pt]{%
          % --> Fichier : image_f95c43.jpg (même dossier que .tex)
          \includegraphics[width=\linewidth, height=5.5cm,
                           keepaspectratio]{image_f95c43}%
        };
      \end{tikzpicture}
    \end{column}

    % Analyse du résultat
    \begin{column}{0.42\textwidth}
      \begin{exampleblock}{\cmark\ Verdict : Cerveau Sain}
        \begin{itemize}\small
          \item Indicateur \textbf{Fiabilité~: Élevée}
          \item Jauge verte animée à \textbf{100\,\%}
          \item Probabilité tumorale \textbf{< seuil~$\tau^*$}
        \end{itemize}
      \end{exampleblock}

      \vspace{6pt}

      \begin{block}{Message affiché}
        \small\itshape
        «~Le modèle a analysé les structures cérébrales
        et n'a détecté aucun signal dépassant le seuil
        de rigueur médicale.~»
      \end{block}

      \vspace{6pt}

      \begin{block}{Choix UX}
        \small
        \textbf{Pas de probabilité brute affichée} — une
        valeur de 79\,\% tumor malgré un verdict sain
        créerait une anxiété injustifiée chez le patient.
      \end{block}

    \end{column}
  \end{columns}

\end{frame}

% ── Slide 4 : Résultat Tumeur ──────────────────────────────
\begin{frame}{Démonstration — Diagnostic~: Tumeur Détectée}

  \begin{columns}[T]

    % Capture d'écran
    \begin{column}{0.55\textwidth}
      \begin{tikzpicture}
        \node[draw=EMSIred!70, line width=1.2pt,
              rounded corners=4pt, inner sep=3pt]{%
          % --> Fichier : image_f96087.jpg (même dossier que .tex)
          \includegraphics[width=\linewidth, height=5.5cm,
                           keepaspectratio]{image_f96087}%
        };
      \end{tikzpicture}
    \end{column}

    % Analyse du résultat
    \begin{column}{0.42\textwidth}
      \begin{alertblock}{\xmark\ Verdict : Tumeur Détectée}
        \begin{itemize}\small
          \item Alerte rouge~(\code{\#ff6b7a}) déclenchée
          \item Score de confiance affiché sur la jauge
          \item $p_{\text{tumeur}} \geq \tau^*_{\text{Youden}}$
        \end{itemize}
      \end{alertblock}

      \vspace{6pt}

      \begin{block}{Inversion des classes — Corrigée}
        \small
        Sans correction~: $p_{\text{brute}} \approx 0.05$\\
        Après~$1 - p_{\text{brute}}$~: $p_{\text{tumeur}} \approx 0.95$\\
        \textbf{Verdict correct~: Tumeur détectée}
      \end{block}

      \vspace{6pt}

      \begin{block}{Avertissement médical}
        \small\itshape
        Résultat fourni à titre académique.
        Consultation d'un radiologue indispensable.
      \end{block}

    \end{column}
  \end{columns}

\end{frame}

% ============================================================
%  SECTION 4 — ÉVALUATION ET PERFORMANCES
% ============================================================
\section{Évaluation et Performances}

\begin{frame}{Évaluation du Modèle}

  \begin{columns}[T]

    % ── Colonne gauche : courbes d'entraînement ────────────
    \begin{column}{0.50\textwidth}
      \begin{block}{Courbes d'Entraînement}
        \vspace{4pt}
        \begin{center}
          % --> Fichier généré par evaluation_kaggle.py : training_history.png
          \includegraphics[width=0.92\linewidth,
                           height=3.6cm,
                           keepaspectratio]{training_history}
        \end{center}
        \vspace{2pt}
        {\scriptsize
          \textbf{Bleu}~: Train \quad \textbf{Rouge}~: Validation\\
          EarlyStopping (patience=5) sur \code{val\_loss}\\
          ReduceLROnPlateau (factor=0.5, patience=3)
        }
      \end{block}
    \end{column}

    % ── Colonne droite : matrice de confusion ──────────────
    \begin{column}{0.47\textwidth}
      \begin{block}{Matrice de Confusion — Set de Test}
        \vspace{6pt}
        \begin{center}
        \renewcommand{\arraystretch}{1.5}
        \begin{tabular}{cc|cc}
          & & \multicolumn{2}{c}{\textbf{Prédit}} \\
          & & \textbf{Sain} & \textbf{Tumeur} \\
          \hline
          \multirow{2}{*}{\rotatebox{90}{\textbf{Réel}}}
            & \textbf{Sain}
              & \cellcolor{EMSIgreen!30}\textbf{VN}
              & \cellcolor{EMSIorange!30}FP \\
            & \textbf{Tumeur}
              & \cellcolor{EMSIred!25}FN
              & \cellcolor{EMSIgreen!30}\textbf{VP} \\
        \end{tabular}
        \end{center}
        \vspace{4pt}
        % --> Remplacez les X par vos valeurs réelles
        \begin{center}
        \begin{tabular}{ll}
          \textcolor{EMSIgreen}{\textbf{VP}} = \textit{X}
            & \textcolor{EMSIorange}{FP} = \textit{X} \\
          \textcolor{EMSIred}{FN} = \textit{X}
            & \textcolor{EMSIgreen}{\textbf{VN}} = \textit{X} \\
        \end{tabular}
        \end{center}
        \vspace{4pt}
        {\scriptsize
          \textbf{Sensibilité}~(Recall)~$= \frac{VP}{VP+FN}$\\[2pt]
          \textbf{Spécificité}~$= \frac{VN}{VN+FP}$\\[2pt]
          \textbf{Seuil~$\tau^*$}~: indice de Youden,
          clamped~$[0.25, 0.75]$
        }
      \end{block}
    \end{column}

  \end{columns}

  \vspace{4pt}
  \begin{center}
  \begin{tikzpicture}
    \node[fill=EMSIlight, draw=EMSIblue2, rounded corners=3pt,
          inner sep=5pt, font=\small]{%
      \textbf{Poids classes~:} $w_{\text{sain}}=1.0$,
      $w_{\text{tumeur}}=2.0$ \quad—\quad
      \textbf{Loss~:} Focal ($\gamma$=2, $\alpha$=0.25) + Label Smoothing ($\varepsilon$=0.1)
    };
  \end{tikzpicture}
  \end{center}

\end{frame}

% ============================================================
%  SECTION 5 — CONCLUSION ET PERSPECTIVES
% ============================================================
\section{Conclusion et Perspectives}

\begin{frame}{Bilan et Perspectives}

  \begin{columns}[T]

    % ── Colonne gauche : bilan ─────────────────────────────
    \begin{column}{0.46\textwidth}
      \begin{exampleblock}{Apports du Projet}
        \begin{itemize}\small
          \item[\cmark] Pipeline Z-Score \textbf{dataset-level}
                        (cohérence train/inférence)
          \item[\cmark] Focal Loss + Label Smoothing
                        contre la \textbf{surconfiance}
          \item[\cmark] Seuil Youden \textbf{automatique}
                        sauvegardé dans \code{.json}
          \item[\cmark] Correction \textbf{inversion classes}
                        par $1 - p_{\text{brute}}$
          \item[\cmark] Interface UX médicale~:
                        \textbf{fiabilité} plutôt que\\
                        probabilité brute
        \end{itemize}
      \end{exampleblock}
    \end{column}

    % ── Colonne droite : perspectives ─────────────────────
    \begin{column}{0.50\textwidth}
      \begin{block}{Perspectives d'Évolution}
        \begin{itemize}\small

          \item \textbf{Segmentation U-Net}\\
                {\scriptsize Produire un masque de localisation
                précis de la lésion, en plus du verdict binaire}

          \vspace{3pt}
          \item \textbf{Cartes de chaleur Grad-CAM}\\
                {\scriptsize Visualiser les zones d'attention du
                modèle sur la coupe IRM → explicabilité clinique}

          \vspace{3pt}
          \item \textbf{Multi-modalité IRM}\\
                {\scriptsize Exploiter T1, T1ce, T2, FLAIR
                simultanément (fusion de features)}

          \vspace{3pt}
          \item \textbf{Calibration isotonique}\\
                {\scriptsize Temperature Scaling pour que
                $p=0.7$ signifie réellement 70\,\% de tumeurs}

          \vspace{3pt}
          \item \textbf{Validation clinique}\\
                {\scriptsize Cohorte externe + audit de biais
                démographique + marquage CE classe IIb}

        \end{itemize}
      \end{block}
    \end{column}

  \end{columns}

\end{frame}

% ── DIAPOSITIVE DE REMERCIEMENT ────────────────────────────
\begin{frame}[plain]
  \begin{tikzpicture}[remember picture, overlay]
    % Fond dégradé
    \fill[EMSIblue] (current page.north west)
      rectangle (current page.south east);
    \fill[EMSIblue2, opacity=0.25]
      (current page.center) circle (5cm);
    % Logo
    \node[anchor=north west, inner sep=10pt]
      at (current page.north west)
      {\includegraphics[height=1.0cm]{logo_emsi}};
  \end{tikzpicture}

  \vspace{1.5cm}
  \begin{center}
    {\Huge\bfseries\color{white} Merci pour votre attention}\\[0.8cm]
    {\large\color{white!85!EMSIblue2}
      \textbf{MOUMEN Doae Malak} \quad \& \quad \textbf{SOUALHINE Douaa}\\[0.3cm]
      Encadrant~: Pr.~BENLAKHDAR Said\\[0.5cm]}
    {\normalsize\color{white!70!EMSIblue2}
      Nous sommes à votre disposition pour les questions.}
  \end{center}
\end{frame}

% ============================================================
\end{document}
